\documentclass[../../../include/open-logic-section]{subfiles}

\begin{document}
\olfileid{sth}{card-arithmetic}{simp}

\olsection{تبسيط الجمع والضرب}

يتبين أن جمع الكاردينالات وضربها فيما جاوز المتناهي بالغَا اليسر.
ومأتاه أن الكاردينالات أعداد ترتيبية مخصوصة، فهي مرتبة ترتيبًا حسنًا
ويتهيأ تناولها على نحو معين. غير أن إقامة البرهان على ذلك ليست باليسيرة.
ولنبدأ بتعريف محكم.

\begin{defn}
نسمي الترتيب~$\canonord$ على أزواج الأعداد الترتيبية
\emph{الترتيب المعياري}، ونجعل
$\tuple{\alpha_1, \alpha_2} \canonord
\tuple{\beta_1, \beta_2}$ إذا وفقط إذا تحقق أحد الشروط الآتية:
\begin{enumerate}
	\item $\max(\alpha_1, \alpha_2) < \max(\beta_1, \beta_2)$؛ أو
	\item $\max(\alpha_1, \alpha_2) = \max(\beta_1, \beta_2)$ و
	$\alpha_1 < \beta_1$؛ أو
	\item $\max(\alpha_1, \alpha_2) = \max(\beta_1, \beta_2)$ و
	$\alpha_1 = \beta_1$ و$\alpha_2 < \beta_2$.
\end{enumerate}
\end{defn}

\begin{lem}
$\tuple{\alpha \times \alpha, \canonord}$ ترتيب جيد، لكل عدد ترتيبي
$\alpha$.
\end{lem}

\begin{proof}
ظاهر أن~$\canonord$ تام على~$\alpha \times \alpha$. فلنفرض أنه لا
$\tuple{\alpha_1, \alpha_2}$ أصغر من
$\tuple{\beta_1, \beta_2}$ في~$\canonord$، ولا الثاني أصغر من الأول.
فيلزم حينئذ
$\max(\alpha_1, \alpha_2) = \max(\beta_1, \beta_2)$ و
$\alpha_1 = \beta_1$ و$\alpha_2 = \beta_2$، ومن ثم
$\tuple{\alpha_1, \alpha_2} = \tuple{\beta_1, \beta_2}$.

وأما حسن الترتيب، فلتكن~$X \subseteq \alpha\times\alpha$ غير
خالية. ولما كانت~$\alpha$ عددًا ترتيبيًا، وُجد~$\delta$ هو أصغر عنصر في
$\Setabs{\max(\gamma_1, \gamma_2)}{\tuple{\gamma_1,
\gamma_2} \in X}$. ثم احذف من
$\Setabs{\tuple{\gamma_1,\gamma_2} \in X}{\max(\gamma_1, \gamma_2) =
\delta}$ كل زوج ليس إحداثيه الأول أصغر ما يكون؛ ومن الباقي اختر ذا
الإحداثي الثاني الأصغر. فهذا أصغر عناصر~$X$ بحسب~$\canonord$.
\end{proof}
\noindent
وها هنا ملاحظة يسيرة.

\begin{prop}\ollabel{simplecardproduct}
إذا كان $\cardeq{\alpha}{\beta}$، فإن
$\cardeq{\alpha \times \alpha}{\beta \times \beta}$.
\end{prop}

\begin{proof}
خذ تقابلًا~$f \colon \alpha \to \beta$ يشهد الفرض؛
فتكون الخريطة
$\tuple{\gamma_1, \gamma_2} \mapsto
\tuple{f(\gamma_1), f(\gamma_2)}$ تقابلًا هي الأخرى، وهذا كاف.
% تصحيح برهاني (0588-N1): قُيدت f بأن تكون تقابلًا واستُعملت تقابلية خريطة الضرب.
\end{proof}
\noindent
ونستعمل ما تقدم لإثبات اللمّة الآتية، وهي عمدة هذا الباب.
\begin{lem}\ollabel{alphatimesalpha}
لكل عدد ترتيبي غير متناهٍ $\alpha$،
$\cardeq{\alpha}{\alpha \times \alpha}$.
\end{lem}

\begin{proof}
نبرهن بالخلف. فليكن~$\alpha$ أصغر عدد ترتيبي غير متناهٍ تخلف عنه
هذا الحكم. وقد ثبت في
\olref[sfr][siz][zigzag]{natsquaredenumerable} أن
$\cardeq{\omega}{\omega\times\omega}$، ولذلك
$\omega \in \alpha$. ثم إن~$\alpha$ كاردينال؛ إذ لو لم يكن كذلك
لكان~$\card{\alpha} \in \alpha$، فيعطينا اختيار ذلك العدد أصغر ما يكون
$\cardeq{\card{\alpha}}{\card{\alpha} \times \card{\alpha}}$،
ويعطينا التعريف~$\cardeq{\card{\alpha}}{\alpha}$. فيلزم عندئذ، بمقتضى
\olref{simplecardproduct}، أن
$\cardeq{\alpha}{\alpha\times\alpha}$، وهو خلف.

والآن خذ أي~$\tuple{\gamma_1, \gamma_2} \in \alpha \times \alpha$،
وانظر في القطعة:
\begin{align*}
	\text{Seg}(\gamma_1, \gamma_2) &= \Setabs{\tuple{\delta_1, \delta_2} \in \alpha \times \alpha}{\tuple{\delta_1, \delta_2} \canonord \tuple{\gamma_1, \gamma_2}}
\end{align*}
ضع $\gamma = \max(\gamma_1, \gamma_2)$، ولاحظ أن
$\tuple{\gamma_1, \gamma_2} \canonord \tuple{\gamma+1, \gamma + 1}$.
فإن كان~$\gamma$ متناهيًا فالقطعة متناهية، ومن ثم كانت كارديناليتها
أصغر من الكاردينال غير المتناهي~$\alpha$. وإن كان~$\gamma$ غير
متناهٍ، حصلنا على:
\begin{align*}
	\text{Seg}(\gamma_1, \gamma_2) &
	\precsim ((\gamma \ordplus 1)\ordtimes (\gamma \ordplus 1))\\
	&\approx (\gamma \ordtimes \gamma)
	\text{، بحسب \olref[ord-arithmetic][using-addition]{ordinfinitycharacter} و
	\olref{simplecardproduct}}\\
	&\approx \gamma \text{، بحسب فرض الاستقراء}\\
	& \prec \alpha\text{، لأن $\alpha$ كاردينال.}
\end{align*}
فإذن~$\ordtype{\alpha\times \alpha, \canonord} \leq \alpha$، ومنه
$\cardle{\alpha \times \alpha}{\alpha}$. وأما
$\cardle{\alpha}{\alpha \times \alpha}$ فظاهر، فتتم النتيجة بمبرهنة
شرودر--برنشتاين.
\end{proof}

وبهذا ننتهي إلى الحكم الذي يبسط هذه العمليات.

\begin{thm}\ollabel{cardplustimesmax}
إذا كان $\cardfont{a}, \cardfont{b}$ كاردينالين غير متناهيين، فإن:
\[
	\cardfont{a}
\cardtimes \cardfont{b} = \cardfont{a} \cardplus \cardfont{b} =
\text{max}(\cardfont{a}, \cardfont{b}).
\]
\end{thm}

\begin{proof}
يجوز، من غير إخلال بالعموم، أن نفرض
$\cardfont{a} = \max(\cardfont{a}, \cardfont{b})$. وعندئذ يعطي
\olref{alphatimesalpha} أن
$\cardfont{a}\cardtimes\cardfont{a} = \cardfont{a} \leq \cardfont{a}
\cardplus \cardfont{b} \leq \cardfont{a} \cardplus \cardfont{a} \leq
\cardfont{a} \cardtimes \cardfont{a}$، فيثبت المطلوب.
\end{proof}\noindent
وكذلك إذا كان~$\cardfont{a}$ غير متناهٍ، فإن اتحاد~$\cardfont{a}$
من المجموعات، لا يزيد حجم واحدة منها على~$\cardfont{a}$، لا يزيد حجمه
هو أيضًا على~$\cardfont{a}$.

\begin{prop}\ollabel{kappaunionkappasize}
ليكن $\cardfont{a}$ كاردينالاً غير متناهٍ. ولكل عدد ترتيبي
$\beta \in \cardfont{a}$، لتكن $X_\beta$ مجموعة تحقق
$\card{X_\beta} \leq \cardfont{a}$. عندئذ
$\card{\bigcup_{\beta \in \cardfont{a}} X_\beta} \leq \cardfont{a}$.
\end{prop}

\begin{proof}
لكل~$\beta \in \cardfont{a}$، ثبّت !!a{injection}
$f_\beta \colon X_\beta \to \cardfont{a}$؛ وفي وجه إمكان هذا
التثبيت كلام يأتي في الحاشية.
\footnote{كيف «ثُبِّتت» هذه
الدوال؟ انظر \olref[sth][choice][countablechoice]{sec}.}
ثم عرّف !!a{injection} $g \colon
\bigcup_{\beta \in \cardfont{a}} X_\beta \to \cardfont{a} \times
\cardfont{a}$ بوضع $g(v) = \tuple{\beta, f_\beta(v)}$، حيث
$v \in X_\beta$ ولا يكون $v \in X_\gamma$ لأي $\gamma \in \beta$.
وبحسب~\olref{cardplustimesmax}،
$\bigcup_{\beta \in \cardfont{a}} X_\beta \preceq
\cardfont{a} \times \cardfont{a} \approx \cardfont{a}$.
\end{proof}

\end{document}
